
\documentclass{article}
\usepackage{colortbl}
\usepackage{makecell}
\usepackage{multirow}
\usepackage{supertabular}

\begin{document}

\newcounter{utterance}

\centering \large Interaction Transcript for game `function\_detective', experiment `numbers\_passive\_labeled\_pairs', episode 39 with Qwen3.6{-}35B{-}A3B{-}FP8{-}without{-}reasoning.
\vspace{24pt}

{ \footnotesize  \setcounter{utterance}{1}
\setlength{\tabcolsep}{0pt}
\begin{supertabular}{c@{$\;$}|p{.15\linewidth}@{}p{.15\linewidth}p{.15\linewidth}p{.15\linewidth}p{.15\linewidth}p{.15\linewidth}}
    \# & \multicolumn{2}{c}{Player} && \multicolumn{2}{c}{Game Master} \\
    \hline

    \theutterance \stepcounter{utterance}  
    & & & \multicolumn{4}{p{0.6\linewidth}}{
        \cellcolor[rgb]{0.9,0.9,0.9}{
            \makecell[{{p{\linewidth}}}]{
                \texttt{\tiny{[P1$\langle$GM]}}
                \texttt{I have a hidden Python function for you to figure out.} \\
\\ 
\texttt{Your job is to infer what it does and then submit a correct Python implementation.} \\
\\ 
\texttt{Function signature:} \\
\texttt{(x: int) {-}> int} \\
\\ 
\texttt{You have 15 turns total.} \\
\\ 
\texttt{I reveal candidate input{-}output pairs one at a time and tell you whether each one is correct.} \\
\\ 
\texttt{You must make your best SOLVE attempt no later than the second{-}to{-}last turn.} \\
\texttt{Do not wait until the very end.} \\
\\ 
\texttt{You can use these actions:} \\
\texttt{{-} JSON format described at the end of the prompt} \\
\texttt{{-} SOLVE: ```python ... ```} \\
\\ 
\texttt{Do not add import statements in your guess.} \\
\texttt{Anything you might need is already available, so just write the function body in hypothesis and return the answer.} \\
\\ 
\texttt{Preloaded examples:} \\
\texttt{1. ((1000, 505), True)} \\
\\ 
\texttt{Scoring:} \\
\texttt{{-} correctness is required for any non{-}zero score} \\
\texttt{{-} in interactive modes, solving with fewer information requests is better} \\
\\ 
\texttt{Input formatting reminders:} \\
\texttt{{-} strings must be quoted} \\
\texttt{{-} booleans must be True/False} \\
\texttt{{-} use commas between multiple values} \\
\texttt{{-} do not write any extra action tags in the same message} \\
\\ 
\texttt{Format example:} \\
\texttt{You write:} \\
\texttt{\{'current\_mode': <probing/confirming/don't know; choose whichever best describes your current state in the task>, 'input\_rationale': <brief description of why you chose this hypothesis>, 'hypothesis': <either the exact string "unknown" if you have no guess yet OR your current best guess as runnable Python matching the given signature, with a one{-}line docstring stating what it does. Never write placeholder code, if you cannot commit to a guess, use "unknown". Match the given signature.} \\
\texttt{e.g.} \\
\texttt{def f(x):} \\
\texttt{	"""Returns the absolute value of x."""} \\
\texttt{	return x if x >= 0 else {-}x} \\
\\ 
\texttt{Use "unknown" if you have no guess.>, 'confidence\_score': <score between 1{-}100. 100 means highest confidence in your hypothesis>\}} \\
\\ 
\texttt{Output only one valid json per turn.} \\
            }
        }
    }
     \\ \\

    \theutterance \stepcounter{utterance}  
    & \multicolumn{4}{p{0.6\linewidth}}{
        \cellcolor[rgb]{0.561,0.737,0.561}{
            \makecell[{{p{\linewidth}}}]{
                \texttt{\tiny{[P1$\rangle$GM]}}
                \texttt{\{'current\_mode': 'probing', 'input\_rationale': 'I have one data point: f(1000) = 505. This suggests the function might be related to reversing digits, summing digits, or some arithmetic operation. 1000 reversed is 0001 = 1, which is not 505. Sum of digits of 1000 is 1, not 505. Let me try a few more inputs to narrow down the pattern. I'll try small numbers first to see if there's a simple pattern like f(x) = x // 2 + something, or f(x) = reverse(x) + something.', 'hypothesis': 'unknown', 'confidence\_score': 10\}} \\
            }
        }
    }
    & & \\ \\

    \theutterance \stepcounter{utterance}  
    & & & \multicolumn{4}{p{0.6\linewidth}}{
        \cellcolor[rgb]{0.9,0.9,0.9}{
            \makecell[{{p{\linewidth}}}]{
                \texttt{\tiny{[P1$\langle$GM]}}
                \texttt{Game Violation: ParseError: Response is not valid JSON.} \\
\texttt{Output ONLY:} \\
\texttt{One JSON with confidence\_score, current\_mode, input\_rationale, and hypothesis keys ...} \\
\texttt{OR} \\
\texttt{SOLVE: ```python ... ```} \\
            }
        }
    }
     \\ \\

    \theutterance \stepcounter{utterance}  
    & \multicolumn{4}{p{0.6\linewidth}}{
        \cellcolor[rgb]{0.561,0.737,0.561}{
            \makecell[{{p{\linewidth}}}]{
                \texttt{\tiny{[P1$\rangle$GM]}}
                \texttt{\{'current\_mode': 'probing', 'input\_rationale': 'I have one data point: f(1000) = 505. This suggests the function might be related to reversing digits, summing digits, or some arithmetic operation. 1000 reversed is 0001 = 1, which is not 505. Sum of digits of 1000 is 1, not 505. Let me try a few more inputs to narrow down the pattern. I will try small numbers first to see if there is a simple pattern like f(x) = x // 2 + something, or f(x) = reverse(x) + something.', 'hypothesis': 'unknown', 'confidence\_score': 10\}} \\
            }
        }
    }
    & & \\ \\

    \theutterance \stepcounter{utterance}  
    & & & \multicolumn{4}{p{0.6\linewidth}}{
        \cellcolor[rgb]{0.9,0.9,0.9}{
            \makecell[{{p{\linewidth}}}]{
                \texttt{\tiny{[P1$\langle$GM]}}
                \texttt{Game Violation: ParseError: Response is not valid JSON.} \\
\texttt{Output ONLY:} \\
\texttt{One JSON with confidence\_score, current\_mode, input\_rationale, and hypothesis keys ...} \\
\texttt{OR} \\
\texttt{SOLVE: ```python ... ```} \\
            }
        }
    }
     \\ \\

    \theutterance \stepcounter{utterance}  
    & \multicolumn{4}{p{0.6\linewidth}}{
        \cellcolor[rgb]{0.561,0.737,0.561}{
            \makecell[{{p{\linewidth}}}]{
                \texttt{\tiny{[P1$\rangle$GM]}}
                \texttt{\{'current\_mode': 'probing', 'input\_rationale': 'I have one data point: f(1000) = 505. I need more data points to determine the pattern. I will test small integers to see if there is a simple arithmetic or digit{-}based relationship.', 'hypothesis': 'unknown', 'confidence\_score': 10\}} \\
            }
        }
    }
    & & \\ \\

    \theutterance \stepcounter{utterance}  
    & & & \multicolumn{4}{p{0.6\linewidth}}{
        \cellcolor[rgb]{0.9,0.9,0.9}{
            \makecell[{{p{\linewidth}}}]{
                \texttt{\tiny{[P1$\langle$GM]}}
                \texttt{Game Violation: ParseError: Response is not valid JSON.} \\
\texttt{Output ONLY:} \\
\texttt{One JSON with confidence\_score, current\_mode, input\_rationale, and hypothesis keys ...} \\
\texttt{OR} \\
\texttt{SOLVE: ```python ... ```} \\
            }
        }
    }
     \\ \\

    \theutterance \stepcounter{utterance}  
    & \multicolumn{4}{p{0.6\linewidth}}{
        \cellcolor[rgb]{0.561,0.737,0.561}{
            \makecell[{{p{\linewidth}}}]{
                \texttt{\tiny{[P1$\rangle$GM]}}
                \texttt{\{"current\_mode": "probing", "input\_rationale": "I have one data point: f(1000) = 505. I need more data points to determine the pattern. I will test small integers to see if there is a simple arithmetic or digit{-}based relationship.", "hypothesis": "unknown", "confidence\_score": 10\}} \\
            }
        }
    }
    & & \\ \\

    \theutterance \stepcounter{utterance}  
    & & & \multicolumn{4}{p{0.6\linewidth}}{
        \cellcolor[rgb]{0.9,0.9,0.9}{
            \makecell[{{p{\linewidth}}}]{
                \texttt{\tiny{[P1$\langle$GM]}}
                \texttt{OUTPUT: ((1, 6), False)} \\
            }
        }
    }
     \\ \\

    \theutterance \stepcounter{utterance}  
    & \multicolumn{4}{p{0.6\linewidth}}{
        \cellcolor[rgb]{0.561,0.737,0.561}{
            \makecell[{{p{\linewidth}}}]{
                \texttt{\tiny{[P1$\rangle$GM]}}
                \texttt{\{"current\_mode": "probing", "input\_rationale": "I have two data points: f(1000) = 505 and f(1) = 6 (incorrect). Since f(1) is not 6, simple digit reversal or sum doesn't work directly. Let's try more inputs to find a pattern. I'll try f(2), f(10), f(100) to see if there's a relationship with powers of 10 or simple arithmetic.", "hypothesis": "unknown", "confidence\_score": 10\}} \\
            }
        }
    }
    & & \\ \\

    \theutterance \stepcounter{utterance}  
    & & & \multicolumn{4}{p{0.6\linewidth}}{
        \cellcolor[rgb]{0.9,0.9,0.9}{
            \makecell[{{p{\linewidth}}}]{
                \texttt{\tiny{[P1$\langle$GM]}}
                \texttt{OUTPUT: ((99, 54), True)} \\
            }
        }
    }
     \\ \\

    \theutterance \stepcounter{utterance}  
    & \multicolumn{4}{p{0.6\linewidth}}{
        \cellcolor[rgb]{0.561,0.737,0.561}{
            \makecell[{{p{\linewidth}}}]{
                \texttt{\tiny{[P1$\rangle$GM]}}
                \texttt{\{"current\_mode": "probing", "input\_rationale": "I have two confirmed data points: f(1000) = 505 and f(99) = 54. Let's analyze these. 1000 reversed is 0001 = 1. 505 is not 1. 99 reversed is 99. 54 is not 99. Let's look at digit sums. Sum of digits of 1000 is 1. Sum of digits of 99 is 18. Neither matches. Let's look at f(x) = reverse(x) + something? No. Let's look at f(x) = x // 2 + something? 1000//2 = 500, +5 = 505. 99//2 = 49, +5 = 54. This looks promising! Let's test this hypothesis: f(x) = x // 2 + 5.", "hypothesis": "unknown", "confidence\_score": 60\}} \\
            }
        }
    }
    & & \\ \\

    \theutterance \stepcounter{utterance}  
    & & & \multicolumn{4}{p{0.6\linewidth}}{
        \cellcolor[rgb]{0.9,0.9,0.9}{
            \makecell[{{p{\linewidth}}}]{
                \texttt{\tiny{[P1$\langle$GM]}}
                \texttt{OUTPUT: (({-}2, 5), False)} \\
            }
        }
    }
     \\ \\

    \theutterance \stepcounter{utterance}  
    & \multicolumn{4}{p{0.6\linewidth}}{
        \cellcolor[rgb]{0.561,0.737,0.561}{
            \makecell[{{p{\linewidth}}}]{
                \texttt{\tiny{[P1$\rangle$GM]}}
                \texttt{\{"current\_mode": "probing", "input\_rationale": "The hypothesis f(x) = x // 2 + 5 failed for x={-}2 (expected 5, got {-}2//2 + 5 = {-}1 + 5 = 4, but the output says False, meaning the actual value is not 4? Or does it mean the pair ({-}2, 5) is incorrect? The format is ((input, output), is\_correct). So f({-}2) is NOT 5. My hypothesis gave 4. So f({-}2) != 5. This doesn't rule out f(x) = x//2 + 5 yet, because 4 != 5. But let's check if f({-}2) could be 4. If f(x) = x//2 + 5, then f({-}2) = 4. The system said (({-}2, 5), False), which means f({-}2) is not 5. It doesn't tell me what f({-}2) IS. I need to test f({-}2) directly or find another pattern. Let's try f(0). If f(x) = x//2 + 5, f(0) = 5. Let's test f(0).", "hypothesis": "unknown", "confidence\_score": 60\}} \\
            }
        }
    }
    & & \\ \\

    \theutterance \stepcounter{utterance}  
    & & & \multicolumn{4}{p{0.6\linewidth}}{
        \cellcolor[rgb]{0.9,0.9,0.9}{
            \makecell[{{p{\linewidth}}}]{
                \texttt{\tiny{[P1$\langle$GM]}}
                \texttt{OUTPUT: ((0, 5), True)} \\
            }
        }
    }
     \\ \\

    \theutterance \stepcounter{utterance}  
    & \multicolumn{4}{p{0.6\linewidth}}{
        \cellcolor[rgb]{0.561,0.737,0.561}{
            \makecell[{{p{\linewidth}}}]{
                \texttt{\tiny{[P1$\rangle$GM]}}
                \texttt{\{"current\_mode": "probing", "input\_rationale": "I have three confirmed data points: f(1000) = 505, f(99) = 54, and f(0) = 5. The hypothesis f(x) = x // 2 + 5 fits all three: 1000//2 + 5 = 505, 99//2 + 5 = 49 + 5 = 54, 0//2 + 5 = 5. Let's test this hypothesis with another value, say f(1). If the hypothesis is correct, f(1) should be 1//2 + 5 = 0 + 5 = 5. Previously, I tested f(1) = 6 and it was False, which is consistent with f(1) = 5. Let me verify f(1) = 5.", "hypothesis": "unknown", "confidence\_score": 80\}} \\
            }
        }
    }
    & & \\ \\

    \theutterance \stepcounter{utterance}  
    & & & \multicolumn{4}{p{0.6\linewidth}}{
        \cellcolor[rgb]{0.9,0.9,0.9}{
            \makecell[{{p{\linewidth}}}]{
                \texttt{\tiny{[P1$\langle$GM]}}
                \texttt{OUTPUT: (({-}1, 5), False)} \\
            }
        }
    }
     \\ \\

    \theutterance \stepcounter{utterance}  
    & \multicolumn{4}{p{0.6\linewidth}}{
        \cellcolor[rgb]{0.561,0.737,0.561}{
            \makecell[{{p{\linewidth}}}]{
                \texttt{\tiny{[P1$\rangle$GM]}}
                \texttt{\{"current\_mode": "probing", "input\_rationale": "The hypothesis f(x) = x // 2 + 5 predicts f({-}1) = {-}1 // 2 + 5 = {-}1 + 5 = 4. The system said (({-}1, 5), False), meaning f({-}1) is not 5. This is consistent with f({-}1) = 4. However, I need to be careful with integer division for negative numbers in Python. {-}1 // 2 is {-}1. So f({-}1) = 4. Let's test f({-}1) = 4 to confirm. Also, let's test f(2) = 2//2 + 5 = 6 to further verify.", "hypothesis": "unknown", "confidence\_score": 80\}} \\
            }
        }
    }
    & & \\ \\

    \theutterance \stepcounter{utterance}  
    & & & \multicolumn{4}{p{0.6\linewidth}}{
        \cellcolor[rgb]{0.9,0.9,0.9}{
            \makecell[{{p{\linewidth}}}]{
                \texttt{\tiny{[P1$\langle$GM]}}
                \texttt{OUTPUT: (({-}1000, {-}495), True)} \\
            }
        }
    }
     \\ \\

    \theutterance \stepcounter{utterance}  
    & \multicolumn{4}{p{0.6\linewidth}}{
        \cellcolor[rgb]{0.561,0.737,0.561}{
            \makecell[{{p{\linewidth}}}]{
                \texttt{\tiny{[P1$\rangle$GM]}}
                \texttt{\{"current\_mode": "confirming", "input\_rationale": "The hypothesis f(x) = x // 2 + 5 predicts f({-}1000) = {-}1000 // 2 + 5 = {-}500 + 5 = {-}495. The system confirmed (({-}1000, {-}495), True). This strongly supports the hypothesis. All tested points (1000{-}>505, 99{-}>54, 0{-}>5, {-}1000{-}>{-}495) are consistent with f(x) = x // 2 + 5. I am confident this is the correct function.", "hypothesis": "def f(x):\textbackslash{}n    \textbackslash{}"\textbackslash{}"\textbackslash{}"Returns x divided by 2 (integer division) plus 5.\textbackslash{}"\textbackslash{}"\textbackslash{}"\textbackslash{}n    return x // 2 + 5", "confidence\_score": 100\}} \\
            }
        }
    }
    & & \\ \\

    \theutterance \stepcounter{utterance}  
    & & & \multicolumn{4}{p{0.6\linewidth}}{
        \cellcolor[rgb]{0.9,0.9,0.9}{
            \makecell[{{p{\linewidth}}}]{
                \texttt{\tiny{[P1$\langle$GM]}}
                \texttt{OUTPUT: ((10, 11), False)} \\
            }
        }
    }
     \\ \\

    \theutterance \stepcounter{utterance}  
    & \multicolumn{4}{p{0.6\linewidth}}{
        \cellcolor[rgb]{0.561,0.737,0.561}{
            \makecell[{{p{\linewidth}}}]{
                \texttt{\tiny{[P1$\rangle$GM]}}
                \texttt{\{"current\_mode": "probing", "input\_rationale": "The hypothesis f(x) = x // 2 + 5 predicts f(10) = 10 // 2 + 5 = 5 + 5 = 10. The system said ((10, 11), False), meaning f(10) is not 11. But my hypothesis gives 10, not 11. So f(10) != 11 is consistent with f(10) = 10. However, I need to check if f(10) is actually 10. Let me test f(10) = 10 to confirm. Also, let's re{-}examine the pattern. f(1000)=505, f(99)=54, f(0)=5, f({-}1000)={-}495. All these fit x//2 + 5. But f(10) should be 10. If f(10) is indeed 10, then the hypothesis holds. Let's test f(10) = 10.", "hypothesis": "unknown", "confidence\_score": 80\}} \\
            }
        }
    }
    & & \\ \\

    \theutterance \stepcounter{utterance}  
    & & & \multicolumn{4}{p{0.6\linewidth}}{
        \cellcolor[rgb]{0.9,0.9,0.9}{
            \makecell[{{p{\linewidth}}}]{
                \texttt{\tiny{[P1$\langle$GM]}}
                \texttt{OUTPUT: ((2, 6), True)} \\
            }
        }
    }
     \\ \\

    \theutterance \stepcounter{utterance}  
    & \multicolumn{4}{p{0.6\linewidth}}{
        \cellcolor[rgb]{0.561,0.737,0.561}{
            \makecell[{{p{\linewidth}}}]{
                \texttt{\tiny{[P1$\rangle$GM]}}
                \texttt{\{"current\_mode": "probing", "input\_rationale": "The hypothesis f(x) = x // 2 + 5 predicts f(2) = 2 // 2 + 5 = 1 + 5 = 6. The system confirmed ((2, 6), True). This is consistent. Let me re{-}check f(10). My hypothesis gives f(10) = 10. The system said ((10, 11), False), which means f(10) is not 11. It doesn't say f(10) is 10. So f(10) could be 10. Let me test f(10) = 10 to confirm the hypothesis.", "hypothesis": "unknown", "confidence\_score": 85\}} \\
            }
        }
    }
    & & \\ \\

    \theutterance \stepcounter{utterance}  
    & & & \multicolumn{4}{p{0.6\linewidth}}{
        \cellcolor[rgb]{0.9,0.9,0.9}{
            \makecell[{{p{\linewidth}}}]{
                \texttt{\tiny{[P1$\langle$GM]}}
                \texttt{OUTPUT: (({-}10, 1), False)} \\
            }
        }
    }
     \\ \\

    \theutterance \stepcounter{utterance}  
    & \multicolumn{4}{p{0.6\linewidth}}{
        \cellcolor[rgb]{0.561,0.737,0.561}{
            \makecell[{{p{\linewidth}}}]{
                \texttt{\tiny{[P1$\rangle$GM]}}
                \texttt{\{"current\_mode": "probing", "input\_rationale": "The hypothesis f(x) = x // 2 + 5 predicts f({-}10) = {-}10 // 2 + 5 = {-}5 + 5 = 0. The system said (({-}10, 1), False), meaning f({-}10) is not 1. This is consistent with f({-}10) = 0. However, I need to verify if f({-}10) is indeed 0. Let me test f({-}10) = 0. Also, let's check f(1) again. My hypothesis gives f(1) = 1 // 2 + 5 = 0 + 5 = 5. Previously, I tested f(1) = 6 and it was False, which is consistent. Let me test f(1) = 5 to confirm.", "hypothesis": "unknown", "confidence\_score": 85\}} \\
            }
        }
    }
    & & \\ \\

    \theutterance \stepcounter{utterance}  
    & & & \multicolumn{4}{p{0.6\linewidth}}{
        \cellcolor[rgb]{0.9,0.9,0.9}{
            \makecell[{{p{\linewidth}}}]{
                \texttt{\tiny{[P1$\langle$GM]}}
                \texttt{OUTPUT: (({-}99, {-}45), True)} \\
            }
        }
    }
     \\ \\

    \theutterance \stepcounter{utterance}  
    & \multicolumn{4}{p{0.6\linewidth}}{
        \cellcolor[rgb]{0.561,0.737,0.561}{
            \makecell[{{p{\linewidth}}}]{
                \texttt{\tiny{[P1$\rangle$GM]}}
                \texttt{\{"current\_mode": "confirming", "input\_rationale": "The hypothesis f(x) = x // 2 + 5 predicts f({-}99) = {-}99 // 2 + 5 = {-}50 + 5 = {-}45. The system confirmed (({-}99, {-}45), True). This is consistent. All tested points so far (1000{-}>505, 99{-}>54, 0{-}>5, {-}1000{-}>{-}495, 2{-}>6, {-}99{-}>{-}45) are consistent with f(x) = x // 2 + 5. The negative results for f(10)=11 and f({-}10)=1 are also consistent with f(10)=10 and f({-}10)=0. I am confident this is the correct function.", "hypothesis": "def f(x):\textbackslash{}n    \textbackslash{}"\textbackslash{}"\textbackslash{}"Returns x divided by 2 (integer division) plus 5.\textbackslash{}"\textbackslash{}"\textbackslash{}"\textbackslash{}n    return x // 2 + 5", "confidence\_score": 100\}} \\
            }
        }
    }
    & & \\ \\

    \theutterance \stepcounter{utterance}  
    & & & \multicolumn{4}{p{0.6\linewidth}}{
        \cellcolor[rgb]{0.9,0.9,0.9}{
            \makecell[{{p{\linewidth}}}]{
                \texttt{\tiny{[P1$\langle$GM]}}
                \texttt{OUTPUT: ((31, 21), False)} \\
            }
        }
    }
     \\ \\

    \theutterance \stepcounter{utterance}  
    & \multicolumn{4}{p{0.6\linewidth}}{
        \cellcolor[rgb]{0.561,0.737,0.561}{
            \makecell[{{p{\linewidth}}}]{
                \texttt{\tiny{[P1$\rangle$GM]}}
                \texttt{\{"current\_mode": "probing", "input\_rationale": "The hypothesis f(x) = x // 2 + 5 predicts f(31) = 31 // 2 + 5 = 15 + 5 = 20. The system said ((31, 21), False), meaning f(31) is not 21. This is consistent with f(31) = 20. However, I need to verify if f(31) is indeed 20. Let me test f(31) = 20. Also, let's check if there's a different pattern. Maybe it's not x//2 + 5. Let's look at the differences: f(1000)=505, f(99)=54, f(0)=5, f({-}1000)={-}495, f(2)=6, f({-}99)={-}45. Let's check f(31) = 20. If f(31) is 20, then the hypothesis holds. Let me test f(31) = 20.", "hypothesis": "unknown", "confidence\_score": 85\}} \\
            }
        }
    }
    & & \\ \\

    \theutterance \stepcounter{utterance}  
    & & & \multicolumn{2}{p{0.3\linewidth}}{
        \cellcolor[rgb]{0.9,0.9,0.9}{
            \makecell[{{p{\linewidth}}}]{
                \texttt{\tiny{[GM$|$GM]}}
                \texttt{Game Verdict: LOSS (turn limit reached)} \\
            }
        }
    }
    & & \\ \\

\end{supertabular}
}

\end{document}
